\scanpage{41}
compatibles avec l'action de $\pi_\infty$ et de
$\unsure{\mathfrak S_3|\mathbb G_3}\deleted{\unsure{\Gamma}}
\times\uline{\{1,\tau\}}$
($\tau$ \deleted{\unsure{est}} \inserted{est} la conjugaison complexe),
il suffit, pour déterminer ces 18 flèches, de déterminer deux d'entre elles,
\deleted{telles que} \inserted{p.ex}
\diagram{\[
\begin{tikzcd}[ampersand replacement=\&,column sep=large,row sep=large]
E(\alpha)\arrow[r]\arrow[d,phantom,"\Vert" description]\&
E'_\infty\arrow[r]\arrow[d,phantom,"\Vert" description]\&E_0\\
\alpha_1\&E_{1/2}\&
\end{tikzcd}
\]}
[NB $E(\alpha)\longrightarrow E_0$ est le composé]
et $=$ de connaître \inserted{les} images dans $E_0,E_1$ de l'élément
marqué $\alpha_\infty$ de $E(\alpha)$,
[d'où éléments marqués
$\unsure{\alpha'_\infty|\beta'_\infty}(i)\in E'_i=E(\alpha'(i))$,
et $\unsure{\alpha_\infty|E\alpha_\infty}(i)\in E_i=E(\alpha(i))$,
$\forall i\in\{0,1,\infty\}$].

On trouve
\[
\begin{aligned}
\alpha'_n(\infty)&=
\left(\sqrt[n]{\tfrac12}\exp\frac{-\pi i}{n},
\sqrt[n]{\tfrac12}\exp\frac{\pi i}{n},1\right)\\
&=(\theta_n^{-1}\overline\zeta_n,\theta_n^{-1}\zeta_n,1)
\end{aligned}
\]
\deleted{$\unsure{\alpha'_\infty(\infty)=}$}
\[
\alpha'_\infty(\infty)=
(\theta_\infty\wedge\overline\zeta_\infty,
\theta_\infty\wedge\zeta_\infty,1)
\]
% archive foliation: 40
